\documentclass[11pt]{article}
\usepackage[margin=1in]{geometry}
\usepackage{amsmath,amssymb,amsthm}
\usepackage{array}
\usepackage{parskip}
\usepackage{booktabs}
\usepackage[UTF8,fontset=fandol]{ctex}

\newtheorem{theorem}{Theorem}
\newtheorem{lemma}[theorem]{Lemma}
\newtheorem{corollary}[theorem]{Corollary}
\newtheorem{proposition}[theorem]{Proposition}
\theoremstyle{definition}
\newtheorem{definition}[theorem]{Definition}
\theoremstyle{remark}
\newtheorem{remark}[theorem]{Remark}

\newcommand{\F}{\mathbb{F}}

\title{Disproof of conjecture \texttt{00000001044}}
\author{earthking11}
\date{2026-09-13}

\begin{document}
\maketitle

\section*{Verdict: FALSE}

We disprove conjecture \texttt{00000001044} as stated. The refutation is
unconditional, elementary, and already complete at the smallest non-trivial
witness $(q,n) = (3,2)$: the hypothesis $d = \gcd(n, q^2-1) > 1$ holds and
$(q-1)/d = 1$, yet \emph{every} $a \in \F_3$ gives a value set whose complement
has size $1$. Hence $a = 0$ is not the unique minimiser, and the other $a$ do
not give a strictly larger value set --- they tie it.

\section{The conjecture, quoted}

Quoted verbatim from \texttt{conjectures/00000001044.md}:

\begin{quote}
\textbf{English.} Definition: The value set of the Dickson polynomial
$D_n(x,a)$ is its image over $\F_q$. Conjecture: When
$\gcd(n, q^2-1) = d > 1$ (the non-permutation case), the minimum size of the
complement of the value set is $(q-1)/d$, attained at $a = 0$ (the monomial
$x^n$); every other $a$ gives a strictly larger value set.
\end{quote}

\begin{quote}
\textbf{中文。} 定义：Dickson 多项式 $D_n(x,a)$ 的值集指其在 $\F_q$ 上的像。
猜想：当 $\gcd(n,q^2-1)=d>1$（非置换情形）时，值集补集的最小尺寸为
$(q-1)/d$，且该极值在 $a=0$（单项式 $x^n$）时取得；一切其他 $a$ 给出严格
更大的值集。
\end{quote}

The claim has two clauses:
\begin{enumerate}
\item[(C1)] the minimum complement size equals $(q-1)/d$ and is attained at
  $a=0$;
\item[(C2)] every other $a$ gives a \emph{strictly larger} value set.
\end{enumerate}
Both fail at $(q,n) = (3,2)$; indeed (C1) and (C2) are mutually inconsistent
as soon as any other $a$ ties $a=0$ (Section~\ref{sec:inconsistent}).

\section{Setup}

\begin{definition}[Dickson polynomial]
For a field $F$ and $a \in F$, the Dickson polynomial $D_n(x,a) \in F[x]$ is
defined by
\[
D_0(x,a) = 2, \qquad D_1(x,a) = x, \qquad
D_k(x,a) = x\,D_{k-1}(x,a) - a\,D_{k-2}(x,a) \quad (k \ge 2).
\]
In particular $D_2(x,a) = x^2 - 2a$ and $D_3(x,a) = x^3 - 3ax$.
\end{definition}

\begin{definition}[Value set and complement]
The \emph{value set} of $D_n(\cdot,a)$ over $\F_q$ is
$V_q(n,a) = \{\,D_n(x,a) : x \in \F_q\,\} \subseteq \F_q$, and its
\emph{complement size} is $\mathrm{comp}_q(n,a) = q - |V_q(n,a)|$.
\end{definition}

Fix $n$ and $q$, and let $d = \gcd(n, q^2-1)$. The conjecture concerns the
regime $d > 1$.

\section{Primary refutation: $(q,n) = (3,2)$}

Take $q = 3$ and $n = 2$, so $D_2(x,a) = x^2 - 2a$ over $\F_3$. Then
\[
d = \gcd(n, q^2-1) = \gcd(2, 9-1) = \gcd(2,8) = 2 > 1,
\qquad
\frac{q-1}{d} = \frac{3-1}{2} = 1 .
\]
So the hypothesis of the conjecture holds strictly.

\begin{theorem}\label{thm:primary}
For $(q,n) = (3,2)$, every $a \in \F_3$ gives a value set of size $2$:
\[
V_3(2,0) = \{0,1\}, \qquad V_3(2,1) = \{1,2\}, \qquad V_3(2,2) = \{0,2\},
\]
hence
\[
\mathrm{comp}_3(2,0) = \mathrm{comp}_3(2,1) = \mathrm{comp}_3(2,2) = 1 .
\]
In particular $a=0$ is not the unique minimiser, and neither $a=1$ nor $a=2$
gives a strictly larger value set than $a=0$.
\end{theorem}

\begin{proof}
Direct evaluation over $\F_3$ of $D_2(x,a) = x^2 - 2a$:
\begin{center}
\begin{tabular}{c c c c}
\toprule
$a$ & $x \mapsto D_2(x,a)$ on $x=0,1,2$ & value set & complement \\
\midrule
$0$ & $0,1,1$ & $\{0,1\}$ & $1$ \\
$1$ & $1,2,2$ & $\{1,2\}$ & $1$ \\
$2$ & $2,0,0$ & $\{0,2\}$ & $1$ \\
\bottomrule
\end{tabular}
\end{center}
Each value set has size $2$, so each complement has size $3-2=1$. All three
complement sizes are equal, and the value sets all have the same size, so none
is strictly larger than another.
\end{proof}

\begin{corollary}
Conjecture \texttt{00000001044} is false. At $(q,n) = (3,2)$ the hypothesis
$d = 2 > 1$ holds, the conjectured value $(q-1)/d = 1$ is correct as a
\emph{minimum}, but it is attained by all three choices $a = 0,1,2$, and no
value set is strictly larger than that of $a=0$.
\end{corollary}

\section{The statement is internally inconsistent}
\label{sec:inconsistent}

Independently of any computation, clauses (C1) and (C2) contradict each other.
Write $v(a) = |V_q(n,a)|$ for the value-set size, so that the complement size
is $q - v(a)$. Clause (C1) says that $a=0$ \emph{minimises the complement},
i.e.\ $q - v(0) \le q - v(a)$, equivalently $v(a) \le v(0)$, for every $a$.
Clause (C2) says that every $a \ne 0$ has a \emph{strictly larger value set},
i.e.\ $v(a) > v(0)$. The two requirements $v(a) \le v(0)$ and $v(a) > v(0)$
cannot both hold for any $a \ne 0$. Thus, whenever $\F_q$ contains more than
one element (so that some $a \ne 0$ exists), the conjecture as written is
unsatisfiable: a minimiser of the complement whose other points all have
strictly larger value sets cannot exist. The plausible intended clause (C2)
would have been ``every other $a$ gives a strictly \emph{smaller} value set''
(i.e.\ a strictly larger complement), which is also what the examples below
refute.

At $(q,n) = (3,2)$ the inconsistency is visible concretely: $a=1$ has
$v(1) = v(0) = 2$, violating the strict inequality $v(1) > v(0)$ demanded by
(C2), while also realising the minimum demanded by (C1).

\section{The same tie at $(5,2)$ and $(7,2)$}

The primary witness is not an isolated accident. For $n = 2$ we have
$D_2(x,a) = x^2 - 2a$, and the value set is the set of squares shifted
by $-2a$; when $2$ is invertible in $\F_q$, as $a$ ranges over $\F_q$ the shift
ranges over all of $\F_q$, so every $a$ produces a translate of the same square
set.

\begin{center}
\begin{tabular}{c c c c c}
\toprule
$(q,n)$ & $d = \gcd(n,q^2-1)$ & $(q-1)/d$ & complement sizes $\mathrm{comp}_q(n,a)$ & value-set sizes \\
\midrule
$(5,2)$ & $\gcd(2,24) = 2$ & $2$ & $a=0,1,2,3,4:\; 2,2,2,2,2$ & all $3$ \\
$(7,2)$ & $\gcd(2,48) = 2$ & $3$ & $a=0,\dots,6:\; 3,3,3,3,3,3,3$ & all $4$ \\
\bottomrule
\end{tabular}
\end{center}

In both rows the hypothesis $d > 1$ and the numerical value $(q-1)/d$ are
correct, but the minimum is attained by \emph{every} $a$, not uniquely by
$a=0$, and no $a$ yields a strictly larger value set. Thus clause (C2) fails
across an infinite family of $n=2$ examples, not merely at $q=3$.

\section{The reversal at $(7,3)$: $a=0$ is the worst choice}

The direction of the claim also fails. Take $q = 7$, $n = 3$, so
$D_3(x,a) = x^3 - 3ax$ over $\F_7$. Then
\[
d = \gcd(3, 49-1) = \gcd(3,48) = 3 > 1,
\qquad
\frac{q-1}{d} = \frac{6}{3} = 2 .
\]

\begin{center}
\begin{tabular}{c c c c}
\toprule
$a$ & value set $V_7(3,a)$ & $|V_7(3,a)|$ & complement \\
\midrule
$0$ & $\{0,1,6\}$ (the cubes) & $3$ & $4$ \\
$1$ & $\{0,2,3,4,5\}$ & $5$ & $2$ \\
$2$ & $\{0,2,3,4,5\}$ & $5$ & $2$ \\
$3$ & $\{0,1,3,4,6\}$ & $5$ & $2$ \\
$4$ & $\{0,2,3,4,5\}$ & $5$ & $2$ \\
$5$ & $\{0,1,3,4,6\}$ & $5$ & $2$ \\
$6$ & $\{0,1,3,4,6\}$ & $5$ & $2$ \\
\bottomrule
\end{tabular}
\end{center}

Here $a = 0$ gives the \emph{largest} complement, $4$, which strictly exceeds
the conjectured minimum $(q-1)/d = 2$, while every nonzero $a$ gives the
smaller complement $2$. So $a=0$ maximises the complement instead of
minimising it: the conjecture has the direction exactly backwards at
$(q,n) = (7,3)$.

\section{Further breakages of the formula and of the ``non-permutation'' label}

Two further defects of the statement are worth recording.

\begin{itemize}
\item \textbf{$(q,n) = (5,3)$.} Here $d = \gcd(3,24) = 3 > 1$, but $d$ does
  not divide $q-1 = 4$, so $(q-1)/d = 4/3$ is not an integer and cannot be a
  set size. Moreover $D_3(x,0) = x^3$ permutes $\F_5$ (as
  $\gcd(3, 5-1) = 1$), so $V_5(3,0) = \F_5$ and $\mathrm{comp}_5(3,0) = 0$.
  Thus in the very case the conjecture labels ``non-permutation''
  ($d>1$), the monomial $a=0$ is in fact a permutation polynomial, and the true
  minimum complement is $0$, not $(q-1)/d$.
\item \textbf{$(q,n) = (5,4)$.} Here $d = \gcd(4,24) = 4$ and
  $(q-1)/d = 1$, but the true minimum complement is $2$ (attained at $a=2,3$),
  and no $a \in \F_5$ attains the claimed value $1$ at all:
  $\mathrm{comp}_5(4,a) = 3,3,2,2,3$ for $a=0,\dots,4$, with
  $\mathrm{comp}_5(4,0) = 3$. So even the numerical value $(q-1)/d$ can fail
  outright.
\end{itemize}

\section{The torus-reading caveat}

One might try to rescue the claim by changing the reading of ``value set''
from the image over $\F_q$ to the image on the norm-$1$ torus, i.e.\ the
subgroup of order $q+1$ of $\F_{q^2}^{\times}$, rather than on $\F_q$. Under
that reading the monomial $a=0$ \emph{is} the unique minimiser, so the
``uniqueness at $a=0$'' half of the claim survives in a qualified sense. But
the minimum complement there is $0$, not $(q-1)/d$; the value $(q-1)/d$ is not
what the torus reading produces. Consequently that reading does not save the
stated formula either, and in any case it contradicts the conjecture's own
definition, which explicitly places the value set over $\F_q$. The refutation
above uses the definition as written.

\section{Reproducibility}

The brute-force checks are carried out by \texttt{reproduce.py}, which uses the
Python~3 standard library only:

\begin{quote}\ttfamily
\$ python3 reproduce.py\\
... \textit{(tables for } (3,2), (5,2), (7,2), (7,3), (5,3), (5,4)\textit{)}\\
PASS: all checks verified; conjecture 00000001044 is FALSE.
\end{quote}

It implements the Dickson recursion over the residues of $\F_q$, computes the
complement size for every $a$, and asserts the $(3,2)$ tie, the $(5,2)$ and
$(7,2)$ ties, and the $(7,3)$ reversal; it prints \texttt{PASS} and exits $0$
exactly when every check holds.

The same facts are formalised in the Lean~4 project under \texttt{lean4/},
using core Lean only (no Mathlib) and no \texttt{sorry}. Its main statements
are

\begin{quote}\ttfamily
theorem conjecture\_00000001044\_false : ...\\
theorem q3n2\_tie : tieWithZero 3 2 = true\\
theorem q7n3\_a0\_largest : \(\forall\) a, a < 7 \(\to\) compSize 7 3 a \(\le\) compSize 7 3 0\\
theorem q5n2\_all : \(\forall\) a, a < 5 \(\to\) compSize 5 2 a = 2\\
theorem q7n2\_all : \(\forall\) a, a < 7 \(\to\) compSize 7 2 a = 3
\end{quote}

\noindent where \texttt{D}, \texttt{valueSet} and \texttt{compSize} are the
Dickson recursion, the value set, and its complement size on residues
represented as naturals reduced modulo $q$. \texttt{lake build} compiles the
project, and \texttt{lake env lean Check.lean} reports the axioms of every
theorem: the concrete computations depend on no axioms at all (the two
divisibility facts use only \texttt{propext}), and none reports
\texttt{sorryAx}. See \texttt{lean4/README.md} for the full statement table.

\section{Files}

\begin{tabular}{l l}
\toprule
File & Description \\
\midrule
\texttt{README.md} & Verdict, quoted conjecture, tables, caveats, file list, reproduction, status. \\
\texttt{main.tex} & This LaTeX source (standalone \texttt{article}, \texttt{tectonic main.tex}). \\
\texttt{build/main.pdf} & PDF produced by \texttt{tectonic}. \\
\texttt{reproduce.py} & Python 3 standard-library brute force; \texttt{PASS}/\texttt{FAIL}, non-zero exit on failure. \\
\texttt{lean4/Main.lean} & Lean 4 formalisation, core Lean only. \\
\texttt{lean4/Check.lean} & \texttt{\#print axioms} audit. \\
\texttt{lean4/README.md} & Statement table and scope note. \\
\bottomrule
\end{tabular}

\end{document}
